\section{Backend: Despliegue de Servicios e Infraestructura Productiva}
\label{sec:s5_backend}

El despliegue productivo materializa la arquitectura monolítica fijada en el Sprint 0
(Sección~\ref{ssec:s0_c4_contenedores}): un único JAR que sirve la API y la SPA, ejecutado en
contenedores. La Figura~\ref{fig:s5_deploy} muestra la topología de despliegue.\par

\begin{figure}[!ht]
    \centering
    \begin{tikzpicture}[
        node distance=0.8cm and 1.2cm,
        box/.style={draw=colorPrimario, fill=headerTabla, thick, rounded corners=2pt,
            text width=2.7cm, align=center, minimum height=1.0cm, font=\sffamily\scriptsize},
        ext/.style={draw=grisISO, fill=grisTecnico, thick, rounded corners=2pt,
            text width=2.5cm, align=center, minimum height=1.0cm, font=\sffamily\scriptsize},
        rel/.style={-{Latex[length=2mm]}, colorPrimario, font=\sffamily\tiny}
    ]
        \node[ext] (cli) {Cliente (navegador)};
        \node[box, right=of cli] (lb) {Balanceador de carga};
        \node[box, right=of lb] (jar) {Contenedor\\JAR EthosHub};
        \node[ext, below=of jar] (supa) {Supabase (PostgreSQL)};
        \node[ext, below=of lb] (env) {Variables de entorno (secretos)};

        \draw[rel] (cli) -- node[above]{HTTPS} (lb);
        \draw[rel] (lb) -- (jar);
        \draw[rel] (jar) -- node[right]{JDBC pooler} (supa);
        \draw[rel] (env) -- node[left]{inyecta} (jar);
    \end{tikzpicture}
    \caption{Topología de despliegue productivo (monolítico en contenedor).}
    \label{fig:s5_deploy}
\end{figure}

\subsection{Configuración Segura de Variables y Credenciales}
\label{ssec:s5_be_secretos}
La configuración de producción se externaliza por completo en \comando{application-prod.yml}
(ignorado por git) mediante placeholders \comando{\$\{VAR\}}. En despliegue solo se exportan las
variables de entorno descritas en la tabla~\ref{tab:s5_env}.

\begin{table}[!ht]
    \centering
    \renewcommand{\arraystretch}{1.3}
    \fontsize{9}{10.8}\selectfont\sffamily
    \begin{tabularx}{\textwidth}{|l|X|}
        \hline
        \rowcolor{colorPrimario}
        \textcolor{white}{\textbf{Variable}} & \textcolor{white}{\textbf{Descripción}} \\ \hline
        \comando{DATABASE\_URL} & JDBC URL de PostgreSQL (Supabase pooler). \\ \hline
        \rowcolor{grisTecnico}
        \comando{SUPABASE\_JWK\_SET\_URI} & JWKS de Supabase para validar JWT. \\ \hline
        \comando{MAIL\_*} & Credenciales SMTP (usuario, contraseña, remitente). \\ \hline
        \rowcolor{grisTecnico}
        \comando{GEMINI\_API\_KEY} & Clave de Google Gemini para IA. \\ \hline
        \comando{CORS\_ALLOWED\_ORIGINS} & Orígenes permitidos en producción. \\ \hline
    \end{tabularx}
    \caption{Variables de entorno de producción (Sprint 5).}
    \label{tab:s5_env}
\end{table}

\begin{BreakingChange}
Las credenciales presentes en el historial de git deben \textbf{rotarse} antes de exponer el
repositorio. Externalizarlas no las borra del pasado: se recomienda reescribir el historial
(\comando{git filter-repo} / BFG) si el repositorio será público.
\end{BreakingChange}

\subsection{Despliegue en Contenedores y Balanceo de Carga}
\label{ssec:s5_be_contenedores}
El despliegue empaqueta FE+BE en un único JAR (vía \comando{build.sh --prod}) y se ejecuta en
contenedores (\comando{Dockerfile}, \comando{docker-compose.yml}) sobre servidores de producción
o nube, con balanceo de carga configurado. Tomcat ajusta \comando{max-http-form-post-size: -1} y
\comando{max-swallow-size: -1} para el correo masivo de administración.

\begin{lstlisting}[caption={Arranque del artefacto monolítico en producción.}, label={lst:s5_run}]
./build.sh --prod
java -jar target/ethoshub-2.0.0.jar --spring.profiles.active=prod
\end{lstlisting}

\subsection{Reversión y Continuidad del Servicio}
\label{ssec:s5_be_rollback}
El despliegue contempla un procedimiento de reversión (\textit{rollback}) ante incidentes, según
la tabla~\ref{tab:s5_rollback}, apoyado en el versionado de la imagen y en los respaldos de la
base de datos (Sección~\ref{ssec:s5_db_backups}).

\begin{table}[!ht]
    \centering
    \renewcommand{\arraystretch}{1.3}
    \fontsize{9}{10.8}\selectfont\sffamily
    \begin{tabularx}{\textwidth}{|l|X|}
        \hline
        \rowcolor{colorPrimario}
        \textcolor{white}{\textbf{Escenario}} & \textcolor{white}{\textbf{Acción de continuidad}} \\ \hline
        Fallo de arranque del JAR & Revertir a la imagen \comando{ethoshub:1.3.0} previa. \\ \hline
        \rowcolor{grisTecnico}
        Migración defectuosa & Restaurar respaldo y reaplicar migración corregida. \\ \hline
        Saturación de tráfico & Escalar réplicas tras el balanceador de carga. \\ \hline
        \rowcolor{grisTecnico}
        Fuga de credenciales & Rotar claves y redeplegar con nuevas variables. \\ \hline
    \end{tabularx}
    \caption{Procedimientos de reversión y continuidad (Sprint 5).}
    \label{tab:s5_rollback}
\end{table}

\subsection{Diagrama de Secuencia: Despliegue Productivo}
\label{ssec:s5_seq_deploy}
La Figura~\ref{fig:s5_seq} muestra la secuencia de despliegue: empaquetado, contenerización con
variables de entorno y conexión a Supabase.

\begin{figure}[!ht]
    \centering
    \begin{tikzpicture}[font=\sffamily\scriptsize, >=Latex,
        part/.style={draw=colorPrimario, fill=headerTabla, rounded corners=2pt,
            minimum height=0.6cm, text width=2.1cm, align=center},
        msg/.style={-{Latex[length=1.6mm]}, colorPrimario},
        ret/.style={-{Latex[length=1.6mm]}, dashed, grisISO}]
        \node[part] (d)  at (0,0)    {DevOps};
        \node[part] (b)  at (3.6,0)  {build.sh --prod};
        \node[part] (dk) at (7.6,0)  {Docker};
        \node[part] (j)  at (10.8,0) {JAR};
        \node[part] (sb) at (13.4,0) {Supabase};
        \foreach \n in {d,b,dk,j,sb} \draw[grisISO, dashed] (\n.south) -- ++(0,-4.7);
        \draw[msg] (0,-0.9)    -- node[above]{1. --prod} (3.6,-0.9);
        \draw[ret] (3.6,-1.6)  -- node[above]{2. JAR 2.0.0} (0,-1.6);
        \draw[msg] (0,-2.3)    -- node[above]{3. build \& run (env)} (7.6,-2.3);
        \draw[msg] (7.6,-3.0)  -- node[above]{4. arranca} (10.8,-3.0);
        \draw[msg] (10.8,-3.7) -- node[above]{5. JDBC pooler} (13.4,-3.7);
    \end{tikzpicture}
    \caption{Diagrama de secuencia del despliegue productivo monolítico.}
    \label{fig:s5_seq}
\end{figure}

\clearpage
